\begin{mistake}{2026-05-03}{03}

\begin{problemcontent}
设 \(f''(1)\) 存在，且 \(f'(1)\ne 0\)，求
\[
\lim_{x\to 1}\left[\frac{1}{f'(1)(x-1)}-\frac{1}{f(x)-f(1)}\right].
\]
\end{problemcontent}

\begin{solutioncontent}
令 \(h=x-1\)，\(a=f'(1)\)，\(b=f''(1)\)，则 \(h\to 0\)，且 \(a\ne 0\)。

由二阶 Taylor 公式（Peano 余项）得
\[
f(1+h)-f(1)=ah+\frac{b}{2}h^2+o(h^2).
\]
因此
\[
\begin{aligned}
\frac{1}{ah}-\frac{1}{f(1+h)-f(1)}
&=\frac{1}{ah}-\frac{1}{ah+\frac{b}{2}h^2+o(h^2)}\\
&=\frac{ah+\frac{b}{2}h^2+o(h^2)-ah}{ah\left(ah+\frac{b}{2}h^2+o(h^2)\right)}\\
&=\frac{\frac{b}{2}h^2+o(h^2)}{a^2h^2+o(h^2)}.
\end{aligned}
\]
分子分母同除以 \(h^2\)，得
\[
\lim_{h\to 0}\frac{\frac{b}{2}+\frac{o(h^2)}{h^2}}{a^2+\frac{o(h^2)}{h^2}}
=\frac{\frac{b}{2}}{a^2}.
\]
代回 \(a=f'(1)\)，\(b=f''(1)\)，原极限为
\[
\boxed{\frac{f''(1)}{2\left[f'(1)\right]^2}}.
\]
\end{solutioncontent}

\mistakeknowledge{
\begin{itemize}
  \item \textbf{核心公式/定理}：若 \(f''(1)\) 存在，则 \(f(1+h)-f(1)=f'(1)h+\frac12 f''(1)h^2+o(h^2)\)。
  \item \textbf{易错点}：不能只用一阶等价 \(f(x)-f(1)\sim f'(1)(x-1)\)，因为两个分式的一阶主部相减后抵消，极限由二阶项决定。
  \item \textbf{关联知识}：遇到 \(\frac{1}{\text{一阶近似}}-\frac{1}{\text{函数增量}}\) 型差值，通常先通分，再把函数增量展开到抵消后的最低非零阶。
\end{itemize}}

\end{mistake}
